\input{../tex-inputs/latex-leseansicht-vorspann}

\section[Rainer Maria Rilke an Arthur Schnitzler, 19. 1. 1902]{L04543 Rainer Maria Rilke an Arthur Schnitzler, 19. 1. 1902}
\nopagebreak\mylabel{L04543v}
\rehead{ }\normalsize\beginnumbering\briefempfaengerindex{Schnitzler, Arthur@\textsc{Schnitzler, Arthur}!zzzRilke, Rainer Maria@\emph{von Rainer Maria Rilke}!1902-01-192@{19. 1. 1902}|(be}
\toendnotes[C]{\smallbreak\pagebreak[2]}
\correspDesc{Versand  durch Rainer Maria Rilke am 19. 1. 1902 in Westerwede
\newline{}Erhalt  durch Arthur Schnitzler im Zeitraum [20. 1. 1902
                  – 25. 1. 1902?] in Wien}\toendnotes[C]{\smallbreak}
\Standort{DLA, A:Rilke-Archiv Gernsbach/Kopien und Abschriften, HS.2022.81.}
\physDesc{Brief, Fotokopie, 1 Blatt, 2 Seiten, 1539 Zeichen
\newline{}Handschrift: schwarze Tinte, deutsche Kurrent
\newline{}Schnitzler: mit Bleistift Vermerk: »\textsc{Rilke}« 
\newline{}Ordnung: 1) Die Fotokopien wurden zwischen 30. 12. 1930 und
                                    14. 1. 1931 durch Carl Sieber\pwindex{Sieber, Carl 1897 – 1945@\textsc{Sieber, Carl} (1897 – 1945)|pw} veranlasst, bevor die
                                 Originale an Schnitzler
                                 zurückgestellt wurden.  2) Der vorliegende Brief wurde am 22. 8. 1937 als Teil
                                 einer umfangreicheren Schenkung von Heinrich Schnitzler\pwindex{Schnitzler, Heinrich 9.\,8.\,1902 Hinterbrühl – 12.\,7.\,1982 Wien@\textsc{Schnitzler, Heinrich} (9.\,8.\,1902 Hinterbrühl – 12.\,7.\,1982 Wien)|pw} an Sándor Wolf\pwindex{Wolf, Sándor 21.\,12.\,1871 Eisenstadt – 2.\,1.\,1946 Haifa@\textsc{Wolf, Sándor} (21.\,12.\,1871 Eisenstadt – 2.\,1.\,1946 Haifa)|pw} übergeben. Die Existenz keines
                                 Objekts der Schenkung ist derzeit nachweisbar.}
\buchAbdrucke{\weitereDrucke{\emph{Rainer Maria Rilke und Arthur Schnitzler. Ihr Briefwechsel.}. Mit Anmerkungen herausgegeben von Heinrich Schnitzler In: \emph{Wort und Wahrheit}, Jg. 13, Nr. 4, April 1958, S. 282–298, hier S. 291–292.} }\toendnotes[C]{\smallbreak}
\pstart
           {\pb}Weſterwede bei Worpswede\oindex{Westerwede@\textbf{Westerwede}|pw}, am 19. Jan. 1902.\pend
           
\pstart{}Mein lieber Doctor \textsc{Schnitzler},\pend\vspace{0.5em}
\pstart
           alles was Sie thun, (und wie lieb, daſs Sie es verſuchen wollen!) iſt mir recht. Ich
               vertraue Ihnen ganz. Sie können \textsc{Bahr\pwindex{Bahr, Hermann 19.\,7.\,1863 Linz – 15.\,1.\,1934 München@\textsc{Bahr, Hermann} (19.\,7.\,1863 Linz – 15.\,1.\,1934 München)|pw}} oder \textsc{Singer\pwindex{Singer, Isidor 16.\,1.\,1857 Budapest – 8.\,12.\,1927 Wien@\textsc{Singer, Isidor} (16.\,1.\,1857 Budapest – 8.\,12.\,1927 Wien)|pw}} ins Vertrauen ziehen und wen Sie{ }ſonſt wollen. – \textsc{Bahr\pwindex{Bahr, Hermann 19.\,7.\,1863 Linz – 15.\,1.\,1934 München@\textsc{Bahr, Hermann} (19.\,7.\,1863 Linz – 15.\,1.\,1934 München)|pw}} iſt ja wohl,{ }ſeit{ }ſeiner Trennung von der »\textsc{Zeit\orgindex{Zeit. Wiener Wochenschrift [Redaktion]@Die Zeit. Wiener Wochenschrift [Redaktion]|pw}}«, an einer \label{K_L04543-1v}\edtext{großen
                  Tageszeitung}{\lemma{\textnormal{\emph{großen
                  Tageszeitung}}}\Cendnote{\textnormal{Tatsächlich arbeitete Bahr\pwindex{Bahr, Hermann 19.\,7.\,1863 Linz – 15.\,1.\,1934 München@\textsc{Bahr, Hermann} (19.\,7.\,1863 Linz – 15.\,1.\,1934 München)|pwk} seit Oktober 1899 für das
                     \emph{Neue Wiener Tagblatt}\orgindex{Neues Wiener Tagblatt@Neues Wiener Tagblatt|pwk} und die \emph{Österreichische Volks-Zeitung}\orgindex{Österreichische Volks-Zeitung@Österreichische Volks-Zeitung|pwk}, die beide zum \emph{Steyrermühl-Konzern}\orgindex{Steyrermühl@Steyrermühl|pwk} gehörten.}}}\label{K_L04543-1} thätig? Für
               eine richtige journaliſtiſche Beſchäftigung fehlt es mir ja wohl an Manchem, aber
               regelmäßige Berichte, beſonders nach außen, könnte ich in gleichem Maße für Zeitungen
               wie für Zeitſchriften auf dem \label{K_L04543-2v}\edtext{neulich}{\lemma{\textnormal{\emph{neulich}}}\Cendnote{\textnormal{XXXX Auszeichnungsfehler: Dokument L04542 nicht gefunden. }}}\label{K_L04543-2} angedeuteten Gebiete
               übernehmen. Wenn ich Ihnen{ }ſpäter Auskünfte über einzelne Punkte geben{ }ſoll, bin ich
               natürlich jeden Augenblick dazu ebenſo bereit, wie ich, wo es noth tut, –{ }ſchriftlich jede
               Verbindung anknüpfen würde, zu der Sie rathen. Und vielen herzlichen Dank für alles.
               Ich hätte auch \textsc{Wassermann\pwindex{Wassermann, Jakob 10.\,3.\,1873 Fürth – 1.\,1.\,1934 Altaussee@\textsc{Wassermann, Jakob} (10.\,3.\,1873 Fürth – 1.\,1.\,1934 Altaussee)|pw}} von meiner Situation geſchrieben, – der ja ein Zeitgenoſſe meiner \label{K_L04543-3v}\edtext{münchner\oindex{München@\textbf{München}|pw} Tage}{\lemma{\textnormal{\emph{münchner Tage}}}\Cendnote{\textnormal{Rilke\pwindex{Rilke, Rainer Maria 4.\,12.\,1875 Prag – 29.\,12.\,1926 Montreux@\textsc{Rilke, Rainer Maria} (4.\,12.\,1875 Prag – 29.\,12.\,1926 Montreux)|pwk} lebte zwischen Herbst 1896 und Sommer 1897 in München\oindex{München@\textbf{München}|pwk}. Durch Jakob Wassermann\pwindex{Wassermann, Jakob 10.\,3.\,1873 Fürth – 1.\,1.\,1934 Altaussee@\textsc{Wassermann, Jakob} (10.\,3.\,1873 Fürth – 1.\,1.\,1934 Altaussee)|pwk}
                  wurde er hier mit Lou Andreas-Salomé\pwindex{Andreas-Salomé, Lou 12.\,2.\,1861 Sankt Petersburg – 5.\,2.\,1937 Göttingen@\textsc{Andreas-Salomé, Lou} (12.\,2.\,1861 Sankt Petersburg – 5.\,2.\,1937 Göttingen)|pwk}
                  bekannt gemacht.}}}\label{K_L04543-3} iſt – aber es fällt mir{ }ſchwer{\dotsfour}
               Wenn {\pb}Sie meinen, daſs er etwas thun kann, – würde ich es
               trotzdem verſuchen.\pend
           
\pstart
           Es iſt eine bange Zeit.\pend
           
\pstart
           Geſtern habe ich Ihnen »\textsc{Das tägliche Leben\pwindex{Rilke, Rainer Maria 4.\,12.\,1875 Prag – 29.\,12.\,1926 Montreux@\textsc{Rilke, Rainer Maria} (4.\,12.\,1875 Prag – 29.\,12.\,1926 Montreux)!tägliche Leben. Drama in zwei Akten@\strich\emph{Das tägliche Leben. Drama in zwei Akten}|pw}}« geſchickt. Nach dem berliner\oindex{Berlin@\textbf{Berlin}|pw} Lacherfolg\eventindex{Residenz-Theater Berlin@\textbf{Residenz-Theater Berlin}!Uraufführung von Das tägliche Leben, 20.12.1901@Uraufführung von Das tägliche Leben, 20.12.1901|pwv} haben \introOben{}mir\introOben{} Direktor \textsc{Lautenburg\pwindex{Lautenburg, Sigmund 11.\,9.\,1851 Budapest – 21.\,7.\,1918 Marienbad@\textsc{Lautenburg, Sigmund} (11.\,9.\,1851 Budapest – 21.\,7.\,1918 Marienbad)|pw}} und der Regiſſeur Doctor \textsc{Zickel\pwindex{Zickel, Martin 1876 Breslau – 1932@\textsc{Zickel, Martin} (1876 Breslau – 1932)|pw}}{ }ſehr
               liebe Briefe geſchrieben voll Vertrauen und Freude an meiner Arbeit. Und in mir{ }ſelbſt iſt das treue Zu-meinem-Stücke-halten auch nicht einen Augenblick erſchüttert
               worden. Es iſt mir lieb, wie vorher. Und deshalb{ }ſende ich es Ihnen. Und würde mich{ }ſehr freuen, wenn es an Ihnen einen Freund gewänne!\pend
           \pstart Dankbar der Ihre: \spacefill\mbox{Rainer Maria Rilke}\pend{}\selectlanguage{ngerman}\endnumbering\briefempfaengerindex{Schnitzler, Arthur@\textsc{Schnitzler, Arthur}!zzzRilke, Rainer Maria@\emph{von Rainer Maria Rilke}!1902-01-192@{19. 1. 1902}|)be}\mylabel{L04543h}
\begin{anhang}
\end{anhang}\newcommand{\dateiname}{L04543}\newcommand{\titel}{Rainer Maria Rilke an Arthur Schnitzler, 19. 1. 1902}\newcommand{\editorInnen}{Herausgegeben von Jahnke, SelmaMüller, Martin Anton}\input{../tex-inputs/latex-leseansicht-abspann}
